\section{Título y descripción del proyecto y del equipo}

\subsection{Nombre del producto/servicio}
\textbf{MisiOps} es una plataforma web de gestión de finanzas personales. Permite a cada usuario registrar sus ingresos y gastos, organizarlos por categorías con presupuesto propio y seguir en un tablero cuánto lleva gastado y cuánto le queda del presupuesto del mes.

\subsection{Breve explicación del problema y la solución propuesta}
Muchas personas no llevan un control de su dinero: llegan a fin de mes con poco o nada ahorrado, o incluso endeudadas, porque no tienen trazabilidad de en qué gastan. Las aplicaciones bancarias muestran movimientos aislados por cuenta, pero no una visión consolidada ni por categoría, y las hojas de cálculo exigen disciplina y tiempo que el usuario típico no tiene.

MisiOps propone una solución de extremo a extremo compuesta por tres piezas desplegadas con Docker Compose:

\begin{itemize}
    \item \textbf{Backend (FastAPI + PostgreSQL)}: expone una API REST versionada (\texttt{/api/v1}) con autenticación JWT, perfil de usuario, categorías por usuario con presupuesto opcional, transacciones con filtros y paginación, y un endpoint de resumen que calcula totales por periodo presupuestario y por categoría.
    \item \textbf{Frontend (Next.js + TypeScript + Tailwind)}: interfaz con registro e inicio de sesión, formulario de registro manual de movimientos, historial agrupado por día con filtros, gestión de categorías y tablero de resumen alimentado por la API.
    \item \textbf{Base de datos (PostgreSQL 16)}: persiste usuarios, categorías y transacciones con integridad referencial (borrado en cascada de usuario a sus datos; una categoría no se elimina ni cambia de tipo mientras tenga transacciones).
\end{itemize}

El proceso de desarrollo sigue un flujo \emph{Spec-Driven Development} (SDD): cada funcionalidad nace como una especificación en \texttt{specs/}, pasa por plan, checklist y tareas, y recién entonces se implementa. Un workflow de GitHub Actions (\texttt{sdd-enforcer}) bloquea cualquier Pull Request con código que no traiga esos artefactos.

\subsection{Integrantes y roles}
\begin{itemize}
    \item \textbf{Jeffrey Monja Castro}: Arquitecto y responsable de DevOps. Definió la estructura del monorepo, la constitución del proyecto y el flujo SDD; construyó los \texttt{Dockerfile}, el \texttt{docker-compose.yaml}, los workflows de CI y las plantillas de Pull Request.
    \item \textbf{Gabriel Espinoza}: Desarrollador Backend. Implementó la API en FastAPI (autenticación, usuarios, categorías y transacciones), los modelos SQLAlchemy asíncronos, y la suite de pruebas con \texttt{pytest} (119 pruebas).
    \item \textbf{Yared Riveros Rodriguez}: Desarrollador Frontend. Diseñó e implementó la aplicación Next.js (pantallas de movimientos, categorías, tablero y login), su integración con la API y las pruebas unitarias del cliente con \texttt{node:test} (39 pruebas).
\end{itemize}
